\section*{Hoofdstuk \ref{ch:g10:concentration-and-dilution} --- Concentratie en verdunning}

\begin{solution}{exo:g10:concentration-and-dilution:1}
$C_m = \qty{5.0}{g} / \qty{0.250}{L} = \qty{20}{g/L}$.
\end{solution}

\begin{solution}{exo:g10:concentration-and-dilution:2}
$c = \qty{0.050}{mol} / \qty{0.200}{L} = \qty{0.25}{mol/L}$.
\end{solution}

\begin{solution}{exo:g10:concentration-and-dilution:3}
$M(\ce{CuSO4}) = 63.5 + 32.1 + 4 \times 16.0 = \qty{159.6}{g/mol}$;
$n = 0.10 \times 0.1000 = \qty{0.0100}{mol}$; $m = 0.0100 \times 159.6 =
\qty{1.60}{g}$.
\end{solution}

\begin{solution}{exo:g10:concentration-and-dilution:4}
$F = 1.0 / 0.050 = 20$; $V_0 = \qty{100.0}{mL} / 20 = \qty{5.0}{mL}$.
\end{solution}

\begin{solution}{exo:g10:concentration-and-dilution:5}
$C_m = c \times M = 0.10 \times 58.5 = \qty{5.85}{g/L}$, ongeveer
\qty{5.9}{g/L}.
\end{solution}

\begin{solution}{exo:g10:concentration-and-dilution:6}
$V_0 = c_1 V_1 / c_0 = 0.020 \times 250.0 / 0.50 = \qty{10.0}{mL}$
(\omterm{def:g10:concentration-and-dilution:dilution}{verdunningsfactor} 25). Een volumepipet van \qty{10.0}{mL} en een
maatkolf van \qty{250.0}{mL}.
\end{solution}

\begin{solution}{exo:g10:concentration-and-dilution:7}
De maatkolf heeft één maatstreep, gemaakt voor één volume en nauwkeurig
tot op een klein deel van een procent; de schaalverdeling van een
bekerglas is maar een ruwe aanwijzing. Zo levert ook een volumepipet één
volume heel nauwkeurig af, terwijl een maatcilinder wijder is en minder
precies af te lezen.
\end{solution}

\begin{solution}{exo:g10:concentration-and-dilution:8}
Tussen \qty{0.02}{mol/L} en \qty{0.04}{mol/L}: ongeveer \qty{0.03}{mol/L}.
Voor meer precisie voeg je tussen die twee extra \omterm{def:g10:concentration-and-dilution:standard}{standaardoplossingen}
toe, bijvoorbeeld met 0.025, 0.030 en \qty{0.035}{mol/L}: een fijnere
reeks waar dat nodig is.
\end{solution}

\begin{solution}{exo:g10:concentration-and-dilution:9}
$M(\ce{C6H12O6}) = \qty{180.0}{g/mol}$; $c = 50 / 180.0 =
\qty{0.28}{mol/L}$.
\end{solution}

\begin{solution}{exo:g10:concentration-and-dilution:10}
$M = \qty{342.0}{g/mol}$, dus $n = \qty{1.00}{mol}$ en $c =
\qty{1.00}{mol/L}$. Tien keer verdund: $c = \qty{0.100}{mol/L}$, dus
$C_m = 0.100 \times 342.0 = \qty{34.2}{g/L}$. Een lepel van
\qty{0.020}{L} bevat $34.2 \times 0.020 = \qty{0.68}{g}$ sacharose.
\end{solution}

\begin{solution}{exo:g10:concentration-and-dilution:11}
(a) Dezelfde massa \omterm{def:g6:solutions-and-solubility:solute}{opgeloste stof} in twee keer zo veel volume:
\qty{6.0}{g/L}. (b) Het volume is ongeveer gehalveerd, de massa \omterm{def:g6:solutions-and-solubility:solute}{opgeloste
stof} onveranderd: ongeveer \qty{24}{g/L}.
\end{solution}

\begin{solution}{exo:g10:concentration-and-dilution:12}
\begin{enumerate}
  \item Drie \omterm{def:g10:concentration-and-dilution:dilution}{verdunningen} met 10 zijn een \omterm{def:g10:concentration-and-dilution:dilution}{verdunning} met $10^3$:
        $\qty{0.10}{mol/L} / 1000 = \qty{1.0e-4}{mol/L}$.
  \item Eén \omterm{def:g10:concentration-and-dilution:dilution}{verdunning} met 1000 zou ofwel een piepklein volume
        \omterm{def:g10:concentration-and-dilution:dilution}{stamoplossing} vragen (\qty{0.1}{mL} in \qty{100}{mL}), dat je met
        een pipet niet nauwkeurig kunt afmeten, ofwel een heel grote kolf
        (\qty{1}{mL} in \qty{1}{L}). Drie tienvoudige \omterm{def:g10:concentration-and-dilution:dilution}{verdunningen}
        gebruiken gewone pipetten en kolven.
\end{enumerate}
\end{solution}

\begin{solution}{exo:g10:concentration-and-dilution:13}
\begin{enumerate}
  \item $M(\ce{CaCl2}) = 40.1 + 2 \times 35.5 = \qty{111.1}{g/mol}$;
        $n = 11.1 / 111.1 = \qty{0.100}{mol}$; $c = 0.100 / 0.5000 =
        \qty{0.200}{mol/L}$.
  \item Elke \ce{CaCl2} geeft één \ce{Ca^2+} en twee \ce{Cl-}:
        $[\ce{Ca^2+}] = \qty{0.200}{mol/L}$ en $[\ce{Cl-}] =
        \qty{0.400}{mol/L}$.
\end{enumerate}
\end{solution}

\begin{solution}{exo:g10:concentration-and-dilution:14}
\begin{enumerate}
  \item $r = \qty{0.60}{cm}$, $h = \qty{0.20}{cm}$:
        $V = \pi \times 0.60^2 \times 0.20 = \qty{0.23}{cm^3}$, dus
        \qty{0.23}{mL} water te veel.
  \item Dezelfde \omterm{def:g6:solutions-and-solubility:solute}{opgeloste stof} zit in \qty{100.23}{mL} in plaats van
        \qty{100.00}{mL}: de concentratie is te laag met $0.23 / 100.23 \approx
        \qty{0.23}{\percent}$.
\end{enumerate}
\end{solution}

\begin{solution}{exo:g10:concentration-and-dilution:15}
$n = 0.20 \times 0.100 + 0.10 \times 0.300 = 0.020 + 0.030 =
\qty{0.050}{mol}$ in \qty{0.400}{L}: $c = \qty{0.125}{mol/L}$. Dat is
niet het gewone gemiddelde, \qty{0.15}{mol/L}: er is drie keer zoveel van
de verdunde oplossing, dus het \omterm{def:g6:pure-substances-and-mixtures:mixture}{mengsel} ligt dichter bij
\qty{0.10}{mol/L}. Het is het gemiddelde gewogen met de volumes.
\end{solution}

\begin{solution}{pb:g10:concentration-and-dilution:1}
\textbf{1.} \qty{0.9}{g} per \qty{0.100}{L}: $C_m = \qty{9.0}{g/L}$.

\textbf{2.} \qty{9.0}{g} per \qty{1000}{mL}, dus \qty{9.0}{mg} per
milliliter.

\textbf{3.} $m = 9.0 \times 0.500 = \qty{4.5}{g}$.

\textbf{4.} Ja: ongeveer \qty{9}{g} per liter tegenover zo'n \qty{360}{g}
per kilogram water bij verzadiging, ongeveer veertig keer minder.

\textbf{5.} $9.0 \times 1.0 = \qty{9.0}{g}$ zout.

\textbf{6.} $M(\ce{NaCl}) = 23.0 + 35.5 = \qty{58.5}{g/mol}$.

\textbf{7.} $c = 9.0 / 58.5 = \qty{0.154}{mol/L}$.

\textbf{8.} $n = 0.154 \times 0.500 = \qty{0.0769}{mol}$ (of
$4.5 / 58.5$).

\textbf{9.} Elke \ce{NaCl} geeft één \ce{Na+} en één \ce{Cl-}: beide met
\qty{0.154}{mol/L}.

\textbf{10.} $0.0769 \times \num{6.02e23} = \num{4.63e22}$ natriumionen.

\textbf{11.} $F = 200 / 9.0 = 22.2$.

\textbf{12.} De \qty{4.5}{g} van de zak komt helemaal uit de
geconcentreerde oplossing: $V_0 = \qty{4.5}{g} / \qty{200}{g/L} = \qty{0.0225}{L} =
\qty{22.5}{mL}$.

\textbf{13.} $V_0 = V_1 / F = 500 / 22.2 = \qty{22.5}{mL}$.

\textbf{14.} $c_0 = 200 / 58.5 = \qty{3.42}{mol/L}$;
$c_0 V_0 = 3.42 \times 0.0225 = \qty{0.0769}{mol}$ en
$c_1 V_1 = 0.154 \times 0.500 = \qty{0.0769}{mol}$: gelijk.

\textbf{15.} Ongeveer $500 - 22.5 = \qty{478}{mL}$ steriel water.

\textbf{16.} Een maatpipet van \qty{25}{mL} (of een buret), af te lezen
tot op \qty{0.1}{mL}.

\textbf{17.} In een maatkolf van \qty{500.0}{mL}, aangevuld tot de
maatstreep.

\textbf{18.} $C_m = 200 \times 23.0 / 500 = \qty{9.20}{g/L}$, te
geconcentreerd met $0.20 / 9.0 \approx \qty{2.2}{\percent}$.

\textbf{19.} \qty{22.5}{mL} van de geconcentreerde (``20 \%'') oplossing
voor één zak van \qty{500}{mL}.
\end{solution}
